\chapter[Neue Drucksachen]{Neue Drucksachen}
\quellseite{11}
Zwischenspiele zu den gebräuchlichsten Chorälen der lutherischen Kirche, bearbeitet von J. A. Theiß, Organist an der Bethlehems-Gemeinde zu Milwaukee, Wis. St. Louis, Mo. Concordia Publishing House. 1903. 191 Seiten 12½ × 9½. In Halbfranz mit Goldtitel gebunden. Preis: \$2.50.

Dieses gerade vor Schluß dieser Nummer fertiggestellte Werk muß als ein verdienstliches und unsern Organisten durchaus empfehlenswertes bezeichnet werden. Wir haben nicht alle, aber doch eine ganze Anzahl dieser Zwischenspiele durchgehen können und können zu ihrer Empfehlung erstens sagen, daß sie weder trivial noch effekthascherisch, sondern würdig gehalten, auch nicht über Einen Leisten geschlagen, sondern gar verschieden geartet sind und doch immer dem Charakter des Chorals entsprechen und sich an ihn auf die eine oder andere Weise anlehnen. Und das ist uns die Hauptsache, und wir können nur wünschen, daß solche und ähnlich gehaltene Zwischenspiele in unsern Kirchen gebraucht werden und die leider noch immer an gar manchen Orten vorkommenden eigenen Einfälle und Phantasien sehr unreifer Organisten verdrängen. — Zum andern ist das Werk auch sehr reichhaltig. Es bietet Zwischenspiele zu 159 Choralmelodieen, und zwar für die gebräuchlicheren und häufiger gespielten Choralmelodien mehr und für die seltener gebrauchten weniger. So finden sich zum Beispiel für den Choral: „Wir glauben all an einen Gott“ 18 Zwischenspiele, „Wer weiß, wie nahe mir mein Ende“ hat 17, „O Lamm Gottes unschuldig“ hat 8, „Nun lob, mein Seel, den Herren“ hat 6 Zwischenspiele. Auch ist darauf Rücksicht genommen, daß manche Choräle in verschiedenen Tonarten gespielt werden, dem Charakter der verschiedenen Texte entsprechend oder aus andern Gründen.

\quellseite{12}
So hat zum Beispiel die Melodie: „O Welt, ich muß dich lassen“ 13 Zwischenspiele in G-dur und 17 in F-dur. Es ist also auch für Mannigfaltigkeit gesorgt und man braucht nicht jeden Sonntag dieselben Zwischenspiele zu hören. Endlich ist auch in Betracht gezogen worden, daß nicht alle Organisten gleich geübt sind; manche Zwischenspiele sind darum leicht und einfach, andere etwas schwieriger. — Die Ausstattung des Buches ist musterhaft. L. F.
